% 合成示例（synthetic example），非正式案例结果
% This project deliberately contains 5 classes of citation problems:
%   1) 2 undefined \cite: vaswani2017attentoin (typo), zhang2020residual (no entry)
%   2) 1 @article missing year: he2016deep
%   3) 1 duplicate \label: sec:intro (defined twice)
%   4) 1 unreferenced bib entry: knuth1984texbook
%   5) 3 normal citations: vaswani2017attention, goodfellow2016deep, lecun2015deep
\documentclass{article}
\usepackage{amsmath}
\title{A Mini Survey Fragment on Sequence Modeling}
\author{Example Author}
\date{}
\begin{document}
\maketitle

\section{Introduction}\label{sec:intro}
The Transformer architecture~\cite{vaswani2017attentoin} changed sequence
modeling, and the original paper~\cite{vaswani2017attention} gives full
details. For deep learning fundamentals see
\cite{goodfellow2016deep,lecun2015deep}. Another work~\cite{zhang2020residual}
discusses training stability of very deep networks; the attention computation
is shown in Eq.~\eqref{eq:attn}.

\section{Background}\label{sec:intro}
\begin{equation}\label{eq:attn}
  \mathrm{Attention}(Q,K,V) =
  \mathrm{softmax}\!\left(\frac{QK^{T}}{\sqrt{d_k}}\right)V
\end{equation}
Residual networks~\cite{he2016deep} ease the training of deep models.

\section{Conclusion}
See Section~\ref{sec:intro} for more background.

\bibliographystyle{plain}
\bibliography{refs}
\end{document}
